\documentclass[11pt]{article}

\usepackage[margin=1in]{geometry}
\usepackage{amsmath,amssymb,amsthm}
\usepackage{mathtools}
\usepackage{enumitem}
\usepackage{booktabs}
\usepackage{array,tabularx}
\usepackage{hyperref}
\usepackage{xurl}
\usepackage{microtype}
% BibLaTeX in-file (single-file paper)

\newcommand{\takeaway}[1]{\paragraph{Takeaway.}\emph{#1}}
\newtheorem{proposition}{Proposition}

\title{Calibration Recipes for Anonymous DHT Deployments:\\ Budgets, Horizons, and Alarm Tax}
\author{}
\date{March 24, 2026 (archive v695)}

\begin{document}
\maketitle

\begin{abstract}
To keep the main Anonymous DHT papers referee-tight and non-redundant, we collect the ``numbers-first'' calibration tables and rule-of-thumb conversions that otherwise get repeated:
(i) witness-amplification horizons for state-dependent leakage,
(ii) translating observation cadence into time-to-detection,
(iii) the inevitable alert-oracle tax for closed-view auditing.
We also include pointers to the CPPC manifest interface (AnonDHT~State~3) without reprinting schemas or cards.\cite{state3}
This note is the numbers-first calibration companion beneath the state/control-plane theorem roots and the anonymous-DHT destination-equalization packet: it owns only recipe tables, conservative conversions, and escalation pointers, while theorem statements, evaluator notarization, runtime-conformance receipts, and public release packaging are imported from companion papers rather than re-owned here.\cite{state1,state2,state3,series:mceq,series:eval1,series:opB,series:abom,series:releaseproperty}
\end{abstract}
\paragraph{Scope.}
This is a calibration bridge: ``numbers-first'' recipes and tables that help size witness cadences, horizons, destination-conditioning consequences, and alarm costs without repeating algebra across the core papers. Within the current stack it is the calibration-side companion beneath State~1, State~2, State~3, and MC-EQ, and it deliberately stops before theorem restatement, evaluator bundles, runtime-plan receipts, or public release packaging.\cite{state1,state2,state3,series:mceq,series:eval1,series:opB,series:abom,series:releaseproperty}

\paragraph{Units.} Budgets are published in bits (use $\log_2$ and $2^{\epsilon}$); results stated in nats can be converted by dividing by $\ln 2$ (see Synthesis~8).\cite{series:notation}

\paragraph{Imports.}
We import the threat/window checklist and trace-interface vocabulary.\cite{series:threatwindows,series:traceifaces}
Amplification bounds are imported from the State series (Papers~1--3), and destination-conditioning / equalization inputs are imported from the committee-contact and MC-EQ notes; identification sizing rules come from the endpoint bridge note.\cite{state1,state2,state3,series:mceq,synth5}
The handoff surfaces for evaluator notarization, runtime conformance, and public claim packaging are imported from Eval~1, Operational~B, ABOM/OINL, and Release~A.\cite{series:eval1,series:opB,series:abom,series:releaseproperty}

\paragraph{Exports.}
We export compact calibration objects: (i) quadratic horizon scaling for persistent witness gaps, (ii) cadence-to-horizon translations, (iii) a tabulated alert-oracle ``tax'' lower bound for closed-view auditing, and (iv) explicit escalation pointers telling the reader which theorem/accountant owner to consult next (State~1, State~2, MC-EQ, or State~3) once a quick recipe has been sized.
For the CPPC manifest interface, see AnonDHT~State~3.\cite{state3}

\paragraph{Boundary and earliest sufficient stop.}
This paper is complete once it publishes the declared input quantities, one conservative translation formula or table per quantity, and an explicit pointer to the next theorem/accountant owner that must be consulted (State~1, State~2, MC-EQ, or State~3, depending on the claim surface).
That makes it the stable numbers-first calibration companion rather than another theorem root, evaluator object, or release note. It does not own new theorem proofs, evaluator bundles, runtime plan certificates / PTL / canary receipts, or ABOM/OINL + release-ledger lineage packaging; those belong to the State series, MC-EQ, Eval~1, Operational~B, ABOM/OINL, and Release~A respectively. The maintained worked artifacts currently name \nolinkurl{calibration_recipe_packet} only as owner-map / series-spine route metadata: \nolinkurl{example_line_item_owner_map.json} names the route, but \nolinkurl{example_verifier_report.json} does not emit that packet, \nolinkurl{example_replay_plans.json} does not include \nolinkurl{eval-calibration-recipe-v1}, and \nolinkurl{example_support_bundle_map.json} does not resolve \nolinkurl{bundle://worked-example/calibration-recipe}. Later terse papers should cite those maintained owners only for the route boundary rather than restating a nonexistent carried packet from memory.\cite{state1,state2,state3,series:mceq,series:eval1,series:opB,series:abom,series:releaseproperty,series:workedexample,series:seriesspines}

\paragraph{Artifact policy.}
This archive is source-only; the full CPPC schema and supporting tooling are available on request.\cite{series:artifactpolicy}

\paragraph{Later-terse-paper citation posture.}
When a later short paper only needs to answer \emph{which conservative calibration object was exported for this anonymous-DHT deployment?}, it should stop at the minimal public calibration-recipe card below plus the tiny horizon-and-tax drill.
If the live issue is instead whether an exact public/on-request calibration packet is materialized, stop at the current-cut boundary: this note exports the calibration card, while the maintained worked artifacts provide only owner-map / series-spine route metadata for \nolinkurl{calibration_recipe_packet}. If it is theorem-root-to-review routing across that already-fixed card, widen only to the maintained series-spine bridge. Other widening owners remain the theorem/accountant roots themselves (State~1, State~2, State~3, or MC-EQ), evaluator notarization (Eval~1), runtime conformance (Operational~B), and public claim / release-lineage packaging (ABOM/OINL and Release~A).\cite{state1,state2,state3,series:mceq,series:eval1,series:opB,series:abom,series:releaseproperty,series:workedexample,series:seriesspines}

\paragraph{Evaluation-series four-stop route.}
Within the local evaluation family the honest route is: Eval~1 for the replayable retry-safe evaluator object, Eval~1A only for the optional checkpoint / witness support packet, Eval~2 for the audit/live transfer contract, and this paper only for the numbers-first conversion layer that turns already-declared quantities into cadence-facing horizon or alert-tax sentences.
This note is therefore the fourth stop, not a substitute for any earlier one.
A later terse Anonymous-DHT paper may cite it together with the earlier evaluation notes, but it may not silently reuse calibration ownership to stand in for retry discipline, support hardening, or detectability transfer.

\begin{table}[t]
\centering
\small
\begin{tabularx}{0.97\linewidth}{@{}p{0.31\linewidth}X@{}}
\toprule
\textbf{Public row} & \textbf{Pinned value}\\
\midrule
Imported claim tuple & \texttt{exposure\_nf\_id = enf.committee\_contact.v1}; \texttt{tw\_id = tw.lookup.30d.v1}; calibration object applies to the committee-contact witness family \\
Witness-gap recipe & Imported per-epoch witness gap $\Delta=0.02$ with target total error $\eta=0.05$ under the State~1 amplification rule of thumb \\
Cadence translation & Public observation cadence fixed at $10^3$ epochs/day, so epoch budgets must be translated into wall-clock time before being quoted \\
Amplification output & $T\approx 2\log_2(1/\eta)/\Delta^2 \approx 2.2\times 10^4$ epochs, hence about $22$ days at the declared cadence \\
Alert policy input & Closed-view alert quality target $(\beta,\alpha)=(0.10,0.01)$ for miss probability and false-alarm probability respectively \\
Alert-tax output & Required public alert budget $\varepsilon_{\mathrm{alert}}\gtrsim \log_2((1-\beta)/\alpha)=\log_2(90)\approx 6.49$ bits \\
Escalation pointer & If the live question becomes a theorem packet, move to State~1 / State~2 / State~3; if it becomes destination equalization, move to MC-EQ \\
Outer-owner routing & Notarization, runtime conformance, and release lineage widen into Eval~1, Operational~B, and ABOM/OINL + Release~A \\
\bottomrule
\end{tabularx}
\caption{A minimal public calibration-recipe card. This is the smallest public answer later terse papers can quote directly when the live issue is the conservative horizon / alert-tax sizing object rather than the theorem root, evaluator bundle, runtime receipt, or release packet.}
\label{tab:minimal-calibration-card}
\end{table}

This card is intentionally non-decorative: it fixes one imported witness family, one gap-and-cadence recipe, one alert-quality target, and one escalation route, without silently re-owning theorem proofs, CPPC manifests, evaluator logs, or release receipts. The maintained worked artifacts name \nolinkurl{calibration_recipe_packet} only in the owner map and generated series-spine route; they do not emit a verifier/support packet for the card in the current cut, and theorem-root-to-review routing remains intentionally delegated to the maintained series-spine bridge.\cite{series:workedexample,series:seriesspines}

\paragraph{One tiny same-card / refreshed-wrapper cutover.}
Keep the imported claim tuple, witness-gap rule, cadence line, alert-quality target, and escalation pointer from Table~\ref{tab:minimal-calibration-card} fixed.
A later terse paper may still reuse the same calibration card if it changes only outer carrier rows such as ABOM/OINL packaging, release-lineage tags, surrounding prose, or companion-paper ordering.
But if it changes the imported witness family, the gap line $(\Delta,\eta)$, the public cadence, the alert target $(\beta,\alpha)$, or the neighboring theorem/accountant owner named by the escalation pointer, then it is not the same calibration object anymore: the deployment must emit a successor calibration card and reopen the relevant theorem/accountant owner rather than claiming that the old numbers were merely ``repackaged.''

\begin{proposition}[Calibration-card non-substitution]
A release may cite Table~\ref{tab:minimal-calibration-card} as a conversion receipt only for the exact imported claim tuple, witness family, $(\Delta,\eta)$ row, cadence row, alert-quality row, and escalation pointer printed in the card. The card is not a theorem packet, not a verifier packet, not a support bundle, and not runtime conformance evidence. In the current cut, \nolinkurl{calibration_recipe_packet} is owner-map / series-spine route metadata only: if \nolinkurl{example_line_item_owner_map.json} names that route while \nolinkurl{example_verifier_report.json}, \nolinkurl{example_replay_plans.json}, and \nolinkurl{example_support_bundle_map.json} do not materialize the corresponding packet / replay hook / support bundle, then any public sentence must say that the current cut does not ship a separate \nolinkurl{calibration_recipe_packet} and must not cite the absent packet as carried evidence.
\end{proposition}

\paragraph{One tiny horizon-and-tax drill.}
With $\Delta=0.02$, $\eta=0.05$, and cadence $10^3$ epochs/day,
\[
T \approx \frac{2\log_2(1/\eta)}{\Delta^2}
   \approx \frac{2\log_2(20)}{0.0004}
   \approx 2.16\times 10^4\text{ epochs}
\]
which is about $21.6$ days at the declared cadence. And with alert target $(\beta,\alpha)=(0.10,0.01)$,
\[
\varepsilon_{\mathrm{alert}} \gtrsim \log_2\!\left(\frac{1-\beta}{\alpha}\right)
= \log_2(90) \approx 6.49\text{ bits}.
\]
A later terse paper may therefore quote ``roughly three weeks of exposure at $10^3$/day and a $6.5$-bit alert tax for a 10\% / 1\% alert target'' and move on, unless the live issue is the underlying theorem packet or a different cadence/alert card.

\paragraph{One tiny successor-support sentence.}
If a later public note only refreshes outer packaging, it may honestly say: \emph{The exported calibration card remained unchanged across the wrapper refresh because the imported witness tuple, $(\Delta,\eta)$ line, cadence line, alert-quality line, and escalation pointer were all reissued verbatim.}
That is the whole reuse test for this note; anything more specific belongs to the theorem/accountant owner or to the maintained owner-map / series-spine route.

\paragraph{One tiny worked-route boundary.}
Suppose a later terse paper needs exactly one honest sentence about the already-fixed calibration object. It may say: \emph{For the declared committee-contact witness family with per-epoch gap $0.02$, cadence $10^3$/day, and alert target $(0.10,0.01)$, the exported calibration card reported roughly a $22$-day amplification horizon and a $6.49$-bit alert tax; the maintained worked owner map and series-spine bridge record where that calibration route sits, but the current cut does not ship a separate \nolinkurl{calibration_recipe_packet}.} That sentence is honest because the local conversion row owns the imported tuple, cadence row, alert-quality row, and quoted horizon/alert-tax outputs. In the current maintained artifacts, \nolinkurl{example_line_item_owner_map.json} names \nolinkurl{calibration_recipe_packet}, but \nolinkurl{example_verifier_report.json} does not emit that packet, \nolinkurl{example_replay_plans.json} does not include \nolinkurl{eval-calibration-recipe-v1}, and \nolinkurl{example_support_bundle_map.json} does not resolve \nolinkurl{bundle://worked-example/calibration-recipe}. So later terse papers may stop at Eval~3 for the numbers-first conversion object itself, cite Synthesis~17/Synthesis~55 only for the owner-map / review-route boundary, reopen State~1 / State~2 / State~3 or MC-EQ when the imported theorem/accountant or destination-equalization row becomes live, and widen to Eval~1 / Eval~2 / Release~A only once the live issue has already left the conversion layer.\cite{state1,state2,state3,series:mceq,series:eval1,series:advantagecontracts,series:releaseproperty,series:workedexample,series:seriesspines}

\paragraph{One tiny owner-stable calibration chain.}
When a later terse paper needs one auditable three-row story rather than one isolated calibration sentence, it should keep the imported theorem row, the local conversion row, and the maintained route / review row separate.

\begin{table}[t]
\centering
\small
\begin{tabularx}{0.97\linewidth}{@{}p{0.22\linewidth}p{0.18\linewidth}X@{}}
\toprule
\textbf{Chain row} & \textbf{Owner} & \textbf{Minimal row that may be quoted without widening further}\\
\midrule
Imported witness row & State~2 in this committee-contact example & \texttt{exposure\_nf\_id = enf.committee\_contact.v1}, \texttt{tw\_id = tw.lookup.30d.v1}, and the declared committee-contact witness family remain imported rather than recalibrated here.\\
Local conversion row & Eval~3 & \texttt{Delta = 0.02}, \texttt{eta = 0.05}, cadence $10^3$/day $\Rightarrow$ about $2.16\times 10^4$ epochs ($\approx 21.6$ days); alert target $(0.10,0.01)$ $\Rightarrow$ about $6.49$ bits of public alert tax.\\
Maintained route / review row & Synthesis~17 and Synthesis~55 & The owner-map / series-spine route for \nolinkurl{calibration_recipe_packet} and the theorem-root-to-review route across the card are imported maintenance surfaces, not local Eval~3 rows and not a materialized verifier/support packet in this cut.\\
\bottomrule
\end{tabularx}
\caption{A minimal owner-stable calibration chain for one committee-contact example. The imported witness row, local conversion row, and maintained route / review row remain separately owned even when a later terse paper quotes them together.}
\label{tab:calibration-chain-card}
\end{table}

\paragraph{One tiny owner-chain drill.}
Suppose a later terse Anonymous-DHT note needs to say three things at once: which witness family was calibrated, what quick horizon / alert-tax numbers were exported, and where the owner-map / series-spine review route now lives.
The honest split is: \emph{State~2 fixes the committee-contact witness family through \texttt{enf.committee\_contact.v1} under \texttt{tw.lookup.30d.v1}.} \emph{Eval~3 separately reports that, for $\Delta=0.02$, $\eta=0.05$, cadence $10^3$/day, and alert target $(0.10,0.01)$, the exported calibration card quotes about a $22$-day horizon and a $6.49$-bit alert tax.} \emph{The owner-map / series-spine route across that card remains maintained by Synthesis~17 and Synthesis~55, and the current cut does not ship a separate verifier/support packet for it.}
Changing the witness family reopens State~2; changing the gap / cadence / alert rows reopens this paper; changing the owner-map / series-spine route reopens the maintained worked example or series-spine bridge.\cite{state2,series:workedexample,series:seriesspines}

\paragraph{One tiny paired-route drill with Eval~2.}
Suppose a later terse note needs both an audit-awareness sentence and a planning sentence from the same deployment.
The honest split is: \emph{Eval~2 fixes that the bounded audited metric $0.78$ transfers only as the live interval $[0.74,0.82]$ under the declared detectability budget.} \emph{Eval~3 separately reports that, for the declared cadence $10^3$/day and alert target $(0.10,0.01)$, the deployment also quotes roughly a $22$-day amplification horizon and a $6.49$-bit alert tax.}
The first sentence remains an audit/live transfer claim and therefore belongs to Eval~2.
The second remains a numbers-first conversion and therefore belongs here.
Changing the detectability packet reopens Eval~2; changing the cadence or alert target reopens this paper.

\paragraph{One tiny paired-route drill with MC-EQ.}
For the destination-equalization subcase, suppose a later terse note needs both a quick planning sentence and the theorem/accountant packet underneath it. The honest split is: \emph{Eval~3 reports that the already-declared destination card with $T=32$ corresponds to about eight days only after the deployment also publishes a planning cadence of four observable routing steps per day.} \emph{MC-EQ / CoverSketch remains the owner of the actual destination card $(k=8,\delta_{eq}^{\mathrm{eff}}=0.005,T=32)$ and its $3.55$-bit session export.} The first sentence belongs here because it adds only a cadence-facing conversion; the second belongs to MC-EQ because it fixes the state / kernel / floor / slack / replay packet. This calibration note therefore cannot replace the destination-equalization witness even when both are cited together.\cite{series:mceq}

\section{Series context and what this note owns}
We use the threat/window checklist \cite{series:threatwindows} and the trace-interface vocabulary \cite{series:traceifaces}.
This note owns only ``back-of-the-envelope'' translations that help readers size cadence, windows, and budgets without re-deriving the same algebra in every paper.
Once a deployment needs a theorem-root claim, a destination-conditioning or equalization packet, a notarized evaluator bundle, a runtime conformance receipt, or a public ABOM/OINL + release-ledger packet, the reader should leave this note and consult the companion owner directly rather than expanding the recipe layer. Owner-map / series-spine wiring for one exported calibration-recipe card belongs to the maintained worked example, and theorem-root-to-review routing across that already-fixed card belongs to the maintained series-spine bridge; this cut does not ship a separate calibration verifier/support packet.\cite{state1,state2,state3,series:mceq,series:eval1,series:opB,series:abom,series:releaseproperty,series:workedexample,series:seriesspines}

\section{Recipe 1: witness amplification needs quadratic horizons}
State series Paper~1 \cite{state1} shows that even a tiny, persistent witness bias amplifies under repetition.
A convenient calibration is the Hoeffding-style rule of thumb
\[
T \approx \frac{2\log_2(1/\eta)}{\Delta^2},
\]
where $\Delta$ is the per-epoch witness gap and $\eta$ is a target total error probability.

\begin{table}[h]
\centering
\small
\begin{tabularx}{\linewidth}{@{}rrr>{\raggedright\arraybackslash}X@{}}
\toprule
Per-epoch gap $\Delta$ & Target error $\eta$ & Rule $T$ & Reading \\
\midrule
0.01 & 0.05 & $\approx 8.6\times 10^4$ & about $86$ days at $10^3$/day (or about $21$ hours at $10^5$/day) \\
0.02 & 0.05 & $\approx 2.2\times 10^4$ & about $22$ days at $10^3$/day (or about $5.2$ hours at $10^5$/day) \\
0.05 & 0.05 & $\approx 3.5\times 10^3$ & about $3.5$ days at $10^3$/day (or about $50$ minutes at $10^5$/day) \\
0.10 & 0.05 & $\approx 8.6\times 10^2$ & about $0.86$ days at $10^3$/day (or about $12$ minutes at $10^5$/day) \\
\bottomrule
\end{tabularx}
\caption{Calibration for witness amplification (Paper~1 \cite{state1}). Epoch vocabularies vary across systems (e.g.\ Nym uses minute-scale epochs \cite{nym_cryptoecon_reward_sharing}); the key scaling is the quadratic dependence on $\Delta$.}
\label{tab:calibration}
\end{table}

\takeaway{If an adversary can extract even a 1--2\% per-epoch witness gap and the system runs at high cadence, ``state'' becomes a fast privacy amplifier unless it is refreshed/privatized.}

\section{Recipe 2: budget-to-horizon translation (cadence \texorpdfstring{$\rightarrow$}{->} time)}
If a monitoring/evidence stream yields $T$ independent-ish epochs per window, the same $T$ can correspond to very different wall-clock horizons depending on cadence and sampling.
Table~\ref{tab:budget-horizon} gives an illustrative translation assuming the rule of thumb above with $\eta=0.05$.

\begin{table}[h]
\centering
\begin{tabular}{@{}rrrr@{}}
\toprule
Per-epoch gap $\Delta$ & $T$ (epochs) & Horizon at $10^3$/day & Horizon at $10^5$/day \\
\midrule
0.01 & $8.6\times 10^4$ & $\approx 86$ days & $\approx 20.7$ hours \\
0.02 & $2.2\times 10^4$ & $\approx 22$ days & $\approx 5.2$ hours \\
0.05 & $3.5\times 10^3$ & $\approx 3.5$ days & $\approx 50$ minutes \\
\bottomrule
\end{tabular}
\caption{Budget-to-horizon calibration for witness amplification (Paper~1 \cite{state1}). If an observer sees only a fraction $q$ of epochs, multiply horizons by $1/q$. If state is refreshed every $W$ days, amplification effectively saturates once $T$ exceeds the refresh horizon.}
\label{tab:budget-horizon}
\end{table}

\section{Recipe 3: the alert-oracle tax is logarithmic but not free}
Closed-view auditing (State series Paper~3 \cite{state3}) can require ``alarms'' or public alerts. If an alert has miss probability $\beta$ and false-alarm probability $\alpha$, the paper's alert-oracle bound yields a necessary privacy cost roughly
\[
\varepsilon \gtrsim \log_2\!\left(\frac{1-\beta}{\alpha}\right),
\]
independent of mechanism details.

\begin{table}[h]
\centering
\begin{tabular}{@{}rrr@{}}
\toprule
Miss prob.\ $\beta$ & False-alarm $\alpha$ & Lower bound $\varepsilon \gtrsim \log_2((1-\beta)/\alpha)$ \\
\midrule
0.10 & 0.05  & 4.17 \\
0.10 & 0.01  & 6.49 \\
0.10 & 0.001 & 9.81 \\
0.05 & 0.05  & 4.25 \\
0.05 & 0.01  & 6.57 \\
0.05 & 0.001 & 9.89 \\
0.01 & 0.05  & 4.31 \\
0.01 & 0.01  & 6.63 \\
0.01 & 0.001 & 9.95 \\
\bottomrule
\end{tabular}
\caption{Calibration for the alert-oracle tax (Paper~3 \cite{state3}) in bits, matching the series-wide publication convention. Tight alarms cost nontrivial $\varepsilon$; pushing $\alpha$ down by a factor of $1000$ adds about $\log_2(1000)\approx 9.97$ bits.}
\label{tab:alert-tax-calibration}
\end{table}

\takeaway{Alarm quality is a privacy knob: to get rare false alarms you must usually spend a meaningful privacy budget, even before any ``payload'' disclosures.}

\section{Recipe 4: endpoint sizing via MaxL and odds inflation (pointer)}
For committee-contact / conditioning claims, begin with the MUCC contact-floor packet in State~2; this note only supplies quick conversions around that packet rather than restating its lower bounds.\cite{state2}
For identification and shortlisting conversions (e.g.\ ``what MaxL budget prevents exact identification among $K$ candidates?''), we do not restate algebra here.
Use the endpoint bridge note \cite{synth5}, which relates realized odds inflation, pointwise maximal leakage, and maximal leakage and records the standard $\log_2(\alpha K)$ sizing rule.

For destination-privacy equalization of routing transcripts, MC-EQ/CoverSketch provides an explicit session scaling bound under uniform covers: $\mathcal{L}_{\max}\le T\log_2(1+2k\delta_{eq})$ (auditable and conservative). This note treats that equalization packet as imported and only converts it into fast calibration guidance.\cite{series:mceq}

\section{CPPC schema pointer}
For the CPPC manifest interface (schema excerpt, example cards, and signing/logging guidance), see the closed-view auditing note (AnonDHT~State~3).\cite{state3} The full schema and tooling are distributed out-of-band on request.\cite{series:artifactpolicy}

\section{Limitations and open problems}
These recipes ignore correlation structure, adaptivity, and adversarial selection effects; they are intended for sanity checks, not tight privacy accounting.
Where correlations matter (e.g.\ retries in the Boss Fight), prefer mechanism-specific bounds and then convert to endpoint budgets via \cite{synth5}.
Likewise, once a table entry is going to appear in a notarized evaluation bundle or a public release packet, the deployment must promote the claim into the companion theorem/equalization/certificate stack instead of treating these recipes as self-sufficient evidence.\cite{state1,state2,series:mceq,series:eval1,series:opB,series:abom,series:releaseproperty}

\section{Takeaways}
\begin{itemize}[leftmargin=*]
\item \textbf{Quadratic amplification:} persistent per-epoch witness gaps ($\Delta$) translate to horizons ($T$) on the order of $\log_2(1/\eta)/\Delta^2$ (Recipe~1).
\item \textbf{Cadence is a privacy knob:} the same $T$ epochs can mean hours or months depending on sampling rate; partial visibility scales horizons by $1/q$ (Recipe~2).
\item \textbf{Alarms cost budget:} tightening false-alarm probability pushes a nontrivial alert-oracle lower bound (Recipe~3); plan to spend privacy budget on operational signals, not just payloads.
\item \textbf{Don't re-derive endpoint sizing:} convert mechanism bounds to endpoint identification / shortlisting budgets via the endpoint bridge note (Recipe~4).
\item \textbf{Escalate cleanly:} once a calibration leaves the scratch-pad stage, move to the theorem / equalization / notarization / runtime-conformance / release owner rather than bloating this recipe note.\cite{state1,state2,state3,series:mceq,series:eval1,series:opB,series:abom,series:releaseproperty}
\end{itemize}

\begin{thebibliography}{99}\small

\bibitem{state1}
Anonymous.
\newblock State-Dependent Anonymity in Anonymous DHTs: Witnesses, Amplification, and Stability.
\newblock 2026.
\newblock AnonDHT State series Paper~1 (this archive).

\bibitem{state2}
Anonymous.
\newblock Committee Contact-Set Privacy in Anonymous DHT Lookups: MUCC, Availability Floors, and Abort-Safe Equalization.
\newblock 2026.
\newblock AnonDHT State series Paper~2 (this archive).

\bibitem{state3}
Anonymous.
\newblock Closed-View Auditing for Stateful Anonymity: Odometers, Alarms, and Machine-Checkable CPPCs.
\newblock 2026.
\newblock AnonDHT State series Paper~3 (this archive).

\bibitem{synth5}
Anonymous.
\newblock Endpoint Metrics Bridge for Anonymous Overlays: Odds Inflation, PML, and MaxL.
\newblock 2026.
\newblock Synthesis series note (Synthesis~5) in this archive.

\bibitem{series:traceifaces}
Anonymous.
\newblock Trace Interfaces for Anonymous DHTs: Observation Tiers, Committee-Contact Traces, and Exposure Normal Forms.
\newblock 2026.
\newblock Synthesis series note (Synthesis~3) in this archive.

\bibitem{series:threatwindows}
Anonymous.
\newblock Threat Models and Repetition Windows for Anonymous DHTs.
\newblock 2026.
\newblock Synthesis series note (Synthesis~4) in this archive.

\bibitem{nym_cryptoecon_reward_sharing}
Xinmu Alexis Cao and others.
\newblock Nym Cryptoeconomic Paper (reward sharing and epochs).
\newblock 2021.
\newblock Public Nym document; used only as an example of an epoch/cadence vocabulary.



\bibitem{series:eval1}
Anonymous.
\newblock Notarized Anonymity Certificates Under Retries.
\newblock 2026.
\newblock Evaluation series Paper~1 (this archive).

\bibitem{series:opB}
Anonymous.
\newblock Auditable Trace Privacy for Anonymous DHTs: Runtime Conformance, Plan Pinning, and Epoch Receipts.
\newblock 2026.
\newblock Operational series Paper~B (this archive).

\bibitem{series:abom}
Anonymous.
\newblock ABOM and OINL for Anonymity Claims: Signed Claim Objects and Reader-Facing Labels.
\newblock 2026.
\newblock Anonymity series Paper~C (this archive).

\bibitem{series:releaseproperty}
Anonymous.
\newblock Release Property Ledger Receipts for Anonymity Claims.
\newblock 2026.
\newblock Release \& destination Paper~A (this archive).

\bibitem{series:mceq}
Anonymous.
\newblock MC-EQ and CoverSketch: Destination Privacy via Markov-Chain Equalization in Anonymous DHT Lookups.
\newblock 2026.
\newblock Release \& destination Paper~B (this archive).

\bibitem{series:workedexample}
Anonymous.
\newblock Worked Example: Receipt Interlock and Source-Only Boundaries.
\newblock 2026.
\newblock Synthesis series Paper~17 (this archive).

\bibitem{series:seriesspines}
Anonymous.
\newblock Series Spines: From Mathematics-Era Roots to Review Handoffs.
\newblock 2026.
\newblock Synthesis series Paper~55 (this archive).

\bibitem{series:artifactpolicy}
Anonymous.
\newblock Artifact policy (source-only archive; artifacts available on request).
\newblock Series note, 2026.

\bibitem{series:notation}
Anonymous.
\newblock Notation and Minimal Model for the One-True-Archive.
\newblock Series synthesis note (Synthesis~8), 2026. (This archive.)



% Added in rev0806 citation-closure hygiene pass.
\bibitem{series:advantagecontracts}
Anonymous.
\newblock Advantage Contracts and Stealth Audits for Anonymous-DHT Privacy Claims.
\newblock Evaluation series Paper~2, 2026. (This archive.)

\end{thebibliography}

\end{document}
